%%%%%%%%%%%%%%%%%%%%%%%%%%%%%%%%%%%%%%%%%%%%%%%%%%%%%%%%%%%%%%%%%%%
%%% Documento LaTeX 											  %%%
%%%%%%%%%%%%%%%%%%%%%%%%%%%%%%%%%%%%%%%%%%%%%%%%%%%%%%%%%%%%%%%%%%%
% Título:	Lista de símbolos
% Autor:  Ignacio Moreno Doblas
% Fecha:  2025-10-29
%%%%%%%%%%%%%%%%%%%%%%%%%%%%%%%%%%%%%%%%%%%%%%%%%%%%%%%%%%%%%%%%%%%
% !TEX root = A0.MiTFG.tex

% ================== Lista de símbolos (sin negrita en la etiqueta) ==================
% ---------------- Lista de símbolos en UNA sola línea ----------------
\chapterbeginx{Símbolos}

% ancho de la columna de símbolos = el del símbolo más largo
\newlength{\symbwidth}
\settowidth{\symbwidth}{DLMS/COSEMM}

\begin{description}[
  style=standard,      % MISMA línea (no nextline)
  font=\normalfont,    % sin negrita en la etiqueta
  labelwidth=\symbwidth,
  labelsep=1em,
  leftmargin=!,
  itemsep=2pt
]

% -------- LATINOS (orden alfabético) --------
\item[$\mathrm{BF}_{\text{disp}}$] Factor de ensanchamiento por dispersión cromática; adimensional
\item[$\mathrm{BF}_{\text{turb}}$] Factor de ensanchamiento por turbulencia; adimensional
\item[$\mathrm{BF}_{\text{total}}$] Factor de ensanchamiento total (dispersión $+$ turbulencia); adimensional
\item[$C$] Parámetro de \textit{chirp} inicial del pulso; adimensional
\item[$c$] Velocidad de la luz en el vacío; \si{\metre\per\second}
\item[$C_{n}^{2}$] Parámetro estructural del índice de refracción; \si{\metre\tothe{-2/3}}
\item[$E_{\text{out}}(t)$] Campo eléctrico real en recepción
\item[$F_{\text{out}}(t)$] Envolvente compleja simulada en recepción; $|F_{\text{out}}|$ es su módulo
\item[$\mathrm{FWHM}$] Ancho temporal a mitad de máxima; \si{\pico\second}
\item[$\mathrm{FWHM}_{\text{disp}}$] Ancho debido únicamente a la dispersión; \si{\pico\second}
\item[$\mathrm{FWHM}_{\text{turb}}$] Ancho debido únicamente a la turbulencia; \si{\pico\second}
\item[$h$] Altura o altitud del enlace sobre el nivel del mar; \si{\metre}
\item[$L_{0}$] Escala externa de von Kármán; \si{\metre}
\item[$L_{D3}$] Longitud de dispersión de tercer orden asociada a $\beta_{3}$; típicamente \si{\kilo\metre}
\item[$L_{D2}$] Longitud de dispersión de segundo orden asociada a $\beta_{2}$; típicamente \si{\kilo\metre}
\item[$l_{0}$] Escala interna de von Kármán; \si{\metre}
\item[$n$] Índice de refracción del medio
\item[$P(h)$] Presión atmosférica a la altura $h$; \si{\pascal}
\item[$T(h)$] Temperatura a la altura $h$; \si{\kelvin}
\item[$U_{\text{out}}(t)$] Envolvente compleja en recepción; $|U_{\text{out}}|$ es la envolvente
\item[$v_\mathrm{f}$] Velocidad de fase en el medio considerado; \si{\metre\per\second}
\item[$z$] Distancia de propagación o longitud del enlace; \si{\kilo\metre}

% -------- GRIEGOS (orden alfabético) --------
\item[$\beta$] Constante de propagación del medio; \si{\per\metre}
\item[$\beta_{2}$] Dispersión de velocidad de grupo; \si{\pico\second\squared\per\kilo\metre}
\item[$\beta_{3}$] Dispersión de tercer orden; \si{\pico\second\cubed\per\kilo\metre}
\item[$\lambda$] Longitud de onda; \si{\nano\metre} o \si{\micro\metre}
\item[$\nu$] Frecuencia; \si{\hertz}

\end{description}

\chapterend